% Part: first-order-logic
% Chapter: sequent-calculus
% Section: rules-and-proofs

\documentclass[../../../include/open-logic-section]{subfiles}

\begin{document}

\iftag{FOL}
      {\olfileid{fol}{seq}{rul}}
      {\olfileid{pl}{seq}{rul}}

\olsection{قواعد و \usetoken{P}{derivation}}

در ادامه، فرض کنید $\Gamma, \Delta, \Pi, \Lambda$ نمایندهٔ توالی‌های
متناهیِ !!{sentence}s باشند.

\begin{defn}[دنباله]
\emph{دنباله} عبارتی است به صورت
\[
\Gamma \Sequent \Delta
\]
که در آن $\Gamma$ و $\Delta$ توالی‌های متناهی (و چه‌بسا تهی) از
!!{sentence}sی زبان $\Lang L$ هستند. $\Gamma$ را \emph{مقدم} و
$\Delta$ را \emph{تالی} می‌نامند.
\end{defn}

\begin{explain}
شهود نهفته در پس یک دنباله چنین است: اگر همهٔ !!{sentence}sی مقدم برقرار
باشند، آنگاه دست‌کم یکی از !!{sentence}sی تالی برقرار است. یعنی اگر
$\Gamma = \tuple{!A_1, \dots, !A_m}$ و
$\Delta = \tuple{!B_1, \dots, !B_n}$ باشند، آنگاه
$\Gamma \Sequent \Delta$ برقرار است اگر و تنها اگر
\[
(!A_1 \land \cdots \land !A_m) \lif (!B_1 \lor \cdots \lor
!B_n)
\]
برقرار باشد. دو حالت خاص وجود دارد: حالتی که $\Gamma$ تهی است و حالتی
که $\Delta$~تهی است. هنگامی که $\Gamma$ تهی است، یعنی $m = 0$،
$\quad \Sequent \Delta$ برقرار است اگر و تنها اگر
$!B_1 \lor \dots \lor !B_n$ برقرار باشد. هنگامی که $\Delta$ تهی است،
یعنی $n = 0$، $\Gamma \Sequent \quad$ برقرار است اگر و تنها اگر
$\lnot(!A_1 \land \dots \land !A_m)$ برقرار باشد. می‌گوییم یک دنباله
معتبر است اگر و تنها اگر !!{sentence}‌ی متناظر با آن معتبر باشد.
\end{explain}

اگر $\Gamma$ توالی‌ای از !!{sentence}s باشد، $\Gamma, !A$ را برای حاصل
افزودن $!A$ به انتهای راست~$\Gamma$ می‌نویسیم (و $!A, \Gamma$ را برای
حاصل افزودن $!A$ به انتهای چپ~$\Gamma$). اگر $\Delta$ نیز توالی‌ای از
!!{sentence}s باشد، آنگاه $\Gamma, \Delta$ الحاق آن دو توالی است.

\begin{defn}[دنبالهٔ آغازین]
\emph{دنبالهٔ آغازین} دنباله‌ای است
\iftag{prvFalse,prvTrue}{به یکی از صورت‌های زیر:
  \begin{enumerate}
    \item $!A \Sequent !A$
    \tagitem{prvTrue}{$\quad \Sequent \ltrue$}{}
    \tagitem{prvFalse}{$\lfalse \Sequent \quad$}{}
  \end{enumerate}}
{به صورت $!A \Sequent !A$}، برای هر !!{sentence} $!A$ در زبان.
\end{defn}

!!^{derivation}s در حساب دنباله‌ای درخت‌های خاصی از دنباله‌ها هستند که
در آن‌ها بالاترین دنباله‌ها دنباله‌های آغازین‌اند، و اگر دنباله‌ای زیر
یک یا دو دنبالهٔ دیگر قرار گیرد، باید به‌درستی با یک قاعدهٔ استنتاج از
آن‌ها به دست آید. قواعد $\Log{LK}$ به دو نوع اصلی تقسیم می‌شوند: قواعد
\emph{منطقی} و قواعد \emph{ساختاری}. قواعد منطقی به نام !!{main
  operator} آن !!{sentence}‌ی خوانده می‌شوند که $!A$ و/یا $!B$ را در
دنبالهٔ پایینی دربر دارد. هر قاعده دو نسخه دارد: یکی برای استنتاج
دنباله‌ای که !!{sentence}‌ی حاوی !!{operator} در سمت چپ آن است، و دیگری
برای دنباله‌ای که !!{sentence} در سمت راست آن است.

\end{document}
